\section{An explicit continuous \texorpdfstring{$\PF_4$}{PF4} separator}
\label{sec:separator}

We now prove Theorem~\ref{thm:separator}.  Put
\[
 c=\frac1{64},\qquad
 B(w)=\sum_{k=-3}^{3}b_kw^k,\qquad
 (b_{-3},\ldots,b_3)=(2,10,23,30,23,10,2),
\]
and define
\begin{equation}\label{eq:separator-kernel}
 f_*(x)=e^{-cx^2}B(e^{2cx})
 =\sum_{k=-3}^{3}b_ke^{ck^2}e^{-c(x-k)^2}.
\end{equation}
The Gaussian-mixture expression proves positivity and Schwartz decay;
palindromy of \((b_k)\) proves evenness.

\subsection{Strict \texorpdfstring{$\PF_3$}{PF3}}

For fixed \(x\), give \(K\in\{-3,\ldots,3\}\) probability proportional to
\(b_ke^{2ckx}\), and write \(\kappa_j\) for its cumulants. Differentiation of
the finite log-partition function gives
\begin{equation}\label{eq:separator-cumulants}
 q=2c-(2c)^2\kappa_2,\qquad
 q'=-(2c)^3\kappa_3,\qquad
 q''=-(2c)^4\kappa_4.
\end{equation}
If \(\mu=\mathbb EK\), then
\(0\leq\mathbb E[(K+3)(3-K)]=9-\kappa_2-\mu^2\).
Thus \(0\leq\kappa_2\leq9\), and hence
\begin{equation}\label{eq:separator-q}
 q\geq2c(1-18c)=2c\frac{23}{32}>0.
\end{equation}
Moreover,
\begin{align}
 F_2
 &=q^3-\{qq''-(q')^2\}\notag\\
 &=q^3+q(2c)^4\kappa_4+(2c)^6\kappa_3^2.       \label{eq:separator-f2}
\end{align}
The identity \(\kappa_4=\mathbb E(K-\mathbb EK)^4-3\kappa_2^2\) gives
\(\kappa_4\geq-243\).  Dropping the final nonnegative term in
\eqref{eq:separator-f2} gives
\[
 F_2\geq q\{q^2-3888c^4\}.
\]
At \(c=1/64\), the two coefficients after division by \(c^2\) are
\[
 q^2/c^2\geq\frac{529}{256},\qquad
 3888c^2=\frac{243}{256}.
\]
Thus \(F_2\geq q c^2(143/128)>0\) globally.  The weighted-mean criterion of
\secref{sec:pf3} proves strict \(\PF_3\).

\subsection{Exact order-four density}

Let \(\epsilon=2c=1/32\), \(t=e^{\epsilon x}\), and
\[
 P(t)=t^3B(t)=2+10t+23t^2+30t^3+23t^4+10t^5+2t^6.
\]
For \(E=t\frac{\dd}{\dd t}\), define
\begin{equation}\label{eq:separator-N-recurrence}
 N_1=tP',\qquad N_{j+1}=t(N_j'P-jN_jP').
\end{equation}
The quotient rule gives \(E^j\log P=N_j/P^j\). Therefore, for
\(\ell_j=(\log f_*)^{(j)}\),
\begin{equation}\label{eq:separator-logjets}
 \ell_2=\frac{-\epsilon P^2+\epsilon^2N_2}{P^2},
 \qquad
 \ell_j=\frac{\epsilon^jN_j}{P^j}\quad(3\leq j\leq6).
\end{equation}
Insert these jets into the central-moment determinant
\eqref{eq:C4-def}. Every monomial has derivative weight twelve, so every
term may first be put over the common denominator \(P^{12}\). The resulting
numerator has the exact factor \(P^8\). After cancellation,
\begin{equation}\label{eq:separator-c4}
 \Cfour[f_*](x)=\frac{\widetilde H(t)}{2^{54}P(t)^4}.
\end{equation}
The exact expansion printed in \appref{app:separator} proves that
\(\widetilde H\) is palindromic of degree 24 and that all 25 integer
coefficients are strictly positive. Hence \(\Cfour[f_*]>0\) on \(\R\).

The established inequalities \(q>0\) and \(F_2>0\) also give
\[
 \kappa=2-Q''=\frac{q^3+F_2}{q^3}>0.
\]
The generic transport identity \eqref{eq:central-identity} now has strictly
positive density and crossing weight for \(f_*\).  Thus
\(\partial_\xi\Psi<0\), and Proposition~\ref{prop:equivalence} proves that
\(f_*\) is strict \(\PF_4\).

\subsection{Nonreal Fourier zeros}

With \(\widehat f(z)=\int_\R f(x)e^{-izx}\dd x\), Gaussian translation gives
\begin{equation}\label{eq:separator-fourier}
 \widehat f_*(z)=\sqrt{\frac{\pi}{c}}e^{-z^2/(4c)}
 \sum_{k=-3}^{3}b_ke^{ck^2}e^{-ikz}.
\end{equation}
Put \(w=e^{-iz}\) and \(u=w+w^{-1}\).  The Laurent factor in
\eqref{eq:separator-fourier} vanishes exactly when
\begin{equation}\label{eq:separator-cubic}
 R_c(u)=2e^{9c}u^3+10e^{4c}u^2
 +(23e^c-6e^{9c})u+30-20e^{4c}
\end{equation}
vanishes. At \(c=1/64\), the exact calculation in
\appref{app:separator} gives \(\operatorname{disc}(R_c)<0\). Thus the real
cubic has a nonreal conjugate pair.
For \(|w|=1\), however, \(w+w^{-1}\) is real. The corresponding \(w\)-roots
therefore lie off the unit circle and lift through \(w=e^{-iz}\) to zeros
with \(\operatorname{Im}z\neq0\). The Gaussian factor is zero-free. This
completes the proof of Theorem~\ref{thm:separator}.
